\subsection{Criteria of continuity and equicontinuity}

Let \(\mathbf{E}\) and \(\mathbf{F}\) be topological vector spaces over the same division ring \(\mathbf{K}\) ; for a linear mapping \(f\) of \(\mathbf{E}\) in \(\mathbf{F}\) to be continuous, it is sufficient for it to be continuous at the origin (GT, III, § 2.8, prop. 23). This proposition generalizes as follows:

PROPOSITION 5. --- Let \(E_{i}(1 \leqslant i \leqslant n)\) and F be topological vector spaces on a non-discrete valued field K. In order that a multilinear mapping f of \(\prod_{i=1}^{n} E_{i}\) in F should be conti-

nuous in the product space \(\prod_{i=1}^{n} E_i\) it is sufficient for it to be continuous at \((0, 0, ..., 0)\) .

Let \((a_{1}, a_{2}, \ldots, a_{n})\) be an arbitrary point of \(\prod_{i=1}^{n} \mathrm{E}_{i}\) ; we must show that for every neighbourhood \(W\) of 0 in \(F\) there exist neighbourhoods \(\mathrm{V}_{i}\) of 0 in \(\mathrm{E}_{i}\) ( \(1 \leqslant i \leqslant n\) ) such that the relations \(z_{i} \in V\) imply

\[
f \left(a _ {1} + z _ {1}, a _ {2} + z _ {2}, \dots , a _ {n} + z _ {n}\right) - f \left(a _ {1}, a _ {2}, \dots , a _ {n}\right) \in \mathrm{W}.
\]

Now, we can write

\[
f (a _ {1} + z _ {1}, \dots , a _ {n} + z _ {n}) - f (a _ {1}, \dots , a _ {n}) = \sum_ {\mathrm{H}} u _ {\mathrm{H}}
\]

where H varies over the \(2^{n}-1\) subsets of the set of integers \(\{1,2,\ldots,n\}\) , excluding the set \(\{1,2,\ldots,n\}\) itself, and where \(u_{\mathrm{H}}=f(y_{1},y_{2},\ldots,y_{n})\) , with \(y_{i}=a_{i}\) if \(i\in H\) and \(y_{i}=z_{i}\) if \(i\notin H\) . There exist \(2^{n}-1\) balanced neighbourhoods \(W_{H}\) of 0 in F such that \(\sum_{H}W_{H}\subset W\) ; on the other hand as f is continuous at \((0,0,\ldots,0)\) by hypothesis, there exists in each \(E_{i}\) a neighbourhood \(U_{i}\) of \(0(1\leqslant i\leqslant n)\) such that the n relations \(x_{i}\in U_{i}\) imply that \(f(x_{1},\ldots,x_{n})\in\bigcap_{\mathrm{H}}\mathrm{W}_{\mathrm{H}}\) . As \(U_{i}\) is absorbent, there exists \(\lambda_{i}\neq0\) in K such that \(\lambda_{i}a_{i}\in U_{i}\) . Let \(\lambda\) be an element of K such that \(|\lambda|\geqslant\prod_{i\in H}|\lambda_{i}|^{-1}\) for each subset H; we show that the neighbourhoods \(V_{i}=\lambda^{-n}U_{i}\) , fulfill the required condition. We can write \(u_{\mathrm{H}}=\mu f(x_{1},\ldots,x_{n})\) where \(x_{i}\in U_{i}\) for \(1\leqslant i\leqslant n\) and \(\mu=\lambda^{-np}(\prod_{i\in H}\lambda_{i}^{-1})\) , p being the number of integers of \(\{1,2,\ldots,n\}\) not in H. From the above \(|\mu|\leqslant1\) , hence \(u_{H}\in\mu W_{H}\subset W_{H}\) since \(W_{H}\) is balanced. The proposition is established.

PROPOSITION 6. --- With the same hypotheses on \(E_{i}(1 \leqslant i \leqslant n)\) and on F as in prop. 5, in order that a set E of multilinear maps of \(\prod_{i=1}^{n} E_{i}\) in F be equicontinuous it is sufficient that the set be equicontinuous at \((0, 0, \ldots, 0)\) .

For, in the demonstration of prop. 5 the \(\mathrm{U}_i\) ( \(1 \leqslant i \leqslant n\) ) can be taken such that the relation \(x_i \in \mathrm{U}_i\) ( \(1 \leqslant i \leqslant n\) ) imply \(f(x_1, \ldots, x_n) \in \bigcap_{\mathrm{H}} \mathrm{W}_{\mathrm{H}}\) for every mapping \(f \in \mathcal{E}\) .

